\chapter{Results and Discussions}

This chapter presents the verification and results of the Agentic AI-driven Digital Twin framework. It begins with the automated testing framework using pytest to validate the correctness of the network topology and execution engine. It then details the results of simulating 50 disruption events across four SME markets, analyzes the economic and temporal deltas, presents baseline versus shocked route comparisons, and discusses the performance of the stateful inventory replenishment layer.

\section{Verification via Automated Testing}
To ensure the mathematical validity and structural integrity of the digital twin and routing engine, a suite of 39 automated unit tests was implemented and executed using the \texttt{pytest} framework.

\subsection{Phase 1 Network Model Verification (11 Tests)}
The network topology graph generated by \texttt{build\_network_model.py} was verified against structural requirements:
\begin{itemize}
    \item \textbf{Graph Topology:} Verified that the graph contains exactly 14 nodes and 19 directed edges.
    \item \textbf{Attribute Coverage:} Ensured all nodes contain mandatory metadata (type, tier, region, sensitivity, role) and all edges contain baseline transit times, costs, and friction penalties.
    \item \textbf{Connectivity:} Confirmed that all four SME consumption markets are reachable via paths starting from the Gulf suppliers.
    \item \textbf{Corridor Bypasses:} Verified that the Cape of Good Hope bypass routes are present and correctly connected.
\end{itemize}

\subsection{Phase 4 Execution Engine Verification (28 Tests)}
The execution engine and Dijkstra pathfinding modules were validated for algorithmic correctness:
\begin{itemize}
    \item \textbf{Graph Immutability:} Verified that shock injection (\texttt{apply\_shock()}) does not mutate the baseline graph, ensuring that subsequent simulations start from a clean state.
    \item \textbf{Shock Calculations:} Checked that critical events increase edge weights and cascade decay factors correctly scale down penalties on downstream edges.
    \item \textbf{Dijkstra Pathfinding:} Validated that Dijkstra optimizations minimize either transit time or logistics cost, and that cost optimization always yields a total cost less than or equal to time optimization.
    \item \textbf{Output File Verification:} Ensured that the execution engine successfully generates the reports \texttt{shock\_analysis\_report.csv}, \texttt{optimized\_routes.json}, and \texttt{phase4\_engine\_summary.json} with valid structure.
\end{itemize}

All 39 unit tests compiled and passed, establishing a verified baseline for running simulation experiments.

\section{Friction Shock Simulation Results}
The framework was evaluated by running the execution engine on 50 simulated friction events, representing varying disruption severity levels. With four target markets, a total of 200 routing scenarios were evaluated. 

\subsection{Economic and Temporal Impact Analysis}
Table~\ref{c5:tab_severity} summarizes the average transit cost deltas across different disruption severity levels:

\begin{table}[htb]
\fontsize{10}{12}\selectfont
\centering
\caption{Disruption Severity and Average Cost Deltas}
\label{c5:tab_severity}
\begin{tabular}{|l|c|c|c|}
\hline
\textbf{Severity Level} & \textbf{Number of Events} & \textbf{Average Cost Delta (USD/MT)} & \textbf{Reroutes Triggered} \\
\hline
Critical & 7 & \$397.98 & 4 \\
High & 15 & \$322.70 & 0 \\
Medium & 18 & \$163.27 & 0 \\
Low & 10 & \$62.09 & 0 \\
\hline
\end{tabular}
\end{table}

The simulation identified the Red Sea Corridor (\texttt{red\_sea}) as the most disrupted node in the network, accounting for the highest cumulative delay and triggering 4 automated rerouting actions during critical events.

\subsection{Route and Corridor Shifts}
When a critical disruption is injected at a chokepoint, the routing engine recalculates paths. Table~\ref{c5:tab_routes} compares the baseline routes against the shocked alternative paths for the Pune and Bengaluru SME markets during a critical Red Sea shock:

\begin{table}[htb]
\fontsize{10}{12}\selectfont
\centering
\caption{Baseline vs. Shocked Rerouting Performance}
\label{c5:tab_routes}
\begin{tabular}{|p{2.5cm}|p{4.5cm}|c|p{4.5cm}|c|}
\hline
\textbf{Market} & \textbf{Baseline Path} & \textbf{Baseline Cost} & \textbf{Shocked Path} & \textbf{Shocked Cost} \\
\hline
Surat SME Market & Oman $\to$ Mundra $\to$ Ahmedabad $\to$ Surat & \$994.60 & Oman $\to$ Mundra $\to$ Ahmedabad $\to$ Surat & \$994.60 \\
\hline
Pune SME Market & Saudi Arabia $\to$ Red Sea $\to$ Mundra $\to$ Mumbai $\to$ Pune & \$2,474.60 & Saudi Arabia $\to$ Strait of Hormuz $\to$ Mundra $\to$ Mumbai $\to$ Pune & \$2,790.20 \\
\hline
Bengaluru SME Market & Saudi Arabia $\to$ Red Sea $\to$ Mundra $\to$ Mumbai $\to$ Bengaluru & \$2,754.60 & Saudi Arabia $\to$ Strait of Hormuz $\to$ Mundra $\to$ Mumbai $\to$ Bengaluru & \$3,070.20 \\
\hline
\end{tabular}
\end{table}

As shown in Table~\ref{c5:tab_routes}, the Surat market route is unaffected by the Red Sea shock because it utilizes Oman as a supplier, bypassing the Red Sea Corridor. For Pune and Bengaluru, which rely on Saudi Arabia, the system automatically routes cargo through the Strait of Hormuz rather than the blocked Red Sea. Although this introduces a cost delta, it prevents shipping halts.

\section{Stateful Inventory and Restocking Performance}
The inventory tracking layer evaluated how logistics delays impact market safety stock levels.

\subsection{Stockout Minimization}
Under normal conditions, markets operate on an $(s, Q)$ replenishment policy, maintaining a safety stock buffer of 5 days of cover. During simulations of low, medium, and high disruptions, the delays did not exceed the days of cover, resulting in zero stockouts. 

During critical disruptions, the transit delay exceeded the days of cover. However, the Cost Analyst agent detected the stockout risk and triggered the \texttt{EmergencyRestockTool}. This restocked the warehouses via emergency air-freight (Premium of \$900/MT). While this incurred premium costs, it successfully reduced the final stockout days to zero across all SME markets.

\subsection{Bullwhip Effect Observation}
The reordering events placed by warehouses were monitored to trace demand amplification. The inventory manager recorded a bullwhip ratio exceeding 1.0 (approximating 1.34) at the Mumbai and Ahmedabad warehouses. This variance amplification occurred because warehouses batch orders and add safety buffers, demonstrating how downstream disruption cascades into order spikes upstream.

\section{Summary}
This chapter presented the results of the project. Automated tests validated the network model and execution engine. Disruption simulations of 50 events demonstrated that the system successfully identifies alternative routing paths (such as Strait of Hormuz and Cape of Good Hope bypasses) and triggers emergency restocking to prevent stockouts at SME markets. The next chapter concludes the report and discusses future work.
